%====================================================================
% INFORME DE ANTEPROYECTO — ESTADO Y DESARROLLO DEL PROYECTO
%
% Compilar con LuaLaTeX, dos veces.
%====================================================================
\documentclass[11pt,letterpaper]{article}

% ---------- Matemática y fuentes ----------
\usepackage{amsmath,mathtools}
\usepackage{amsthm}
\usepackage{fontspec}
\usepackage{unicode-math}
\setmainfont{Latin Modern Roman}
\setmathfont{Latin Modern Math}

% ---------- Idioma ----------
\usepackage{polyglossia}
\setmainlanguage{spanish}

% ---------- Diseño ----------
\usepackage[margin=1in]{geometry}
\usepackage{microtype}
\usepackage{enumitem}
\usepackage[table]{xcolor}
\usepackage{booktabs}
\usepackage{listings}
\usepackage{fancyhdr}
\setlength{\headheight}{14pt}
\pagestyle{fancy}
\fancyhf{}
\renewcommand{\headrulewidth}{0pt}
\fancyhead[R]{\small\itshape Informe de avance del proyecto}
\fancyfoot[C]{\small\thepage}

\lstset{
	language=Python, basicstyle=\ttfamily\small, frame=lines,
	keywordstyle=\bfseries, commentstyle=\itshape\color{gray},
	showstringspaces=false, columns=fullflexible, xleftmargin=1em
}

% ---------- Enlaces y referencias cruzadas ----------
\usepackage[colorlinks=true,linkcolor=blue!40!black,citecolor=blue!40!black,urlcolor=blue!40!black]{hyperref}
\usepackage[capitalise]{cleveref}
\crefname{theorem}{Teorema}{Teoremas}
\crefname{proposition}{Proposición}{Proposiciones}
\crefname{lemma}{Lema}{Lemas}
\crefname{corollary}{Corolario}{Corolarios}
\crefname{definition}{Definición}{Definiciones}
\crefname{example}{Ejemplo}{Ejemplos}
\crefname{conjecture}{Conjetura}{Conjeturas}
\crefname{section}{Sección}{Secciones}
\crefname{table}{Tabla}{Tablas}
\crefname{equation}{ecuación}{ecuaciones}

% ---------- Teoremas (numeración por sección) ----------
\theoremstyle{plain}
\newtheorem{theorem}{Teorema}[section]
\newtheorem{proposition}[theorem]{Proposición}
\newtheorem{lemma}[theorem]{Lema}
\newtheorem{corollary}[theorem]{Corolario}
\newtheorem{conjecture}[theorem]{Conjetura}
\theoremstyle{definition}
\newtheorem{definition}[theorem]{Definición}
\newtheorem{example}[theorem]{Ejemplo}
\theoremstyle{remark}
\newtheorem{remark}[theorem]{Observación}
\numberwithin{equation}{section}


% Texto predeterminado (aparece en rojo y entre corchetes)
\newcommand{\ph}[1]{\textcolor{red!70!black}{\textit{[#1]}}}

% Estados para la tabla de avance
\newcommand{\done}{\textcolor{green!45!black}{\textbf{Completado}}}
\newcommand{\wip}{\textcolor{orange!80!black}{\textbf{En curso}}}
\newcommand{\todo}{\textcolor{red!70!black}{\textbf{Pendiente}}}

% =====================================================================
% DATOS DEL DOCUMENTO  (edita aquí)
% =====================================================================
\title{\textbf{\ph{Burgers, Hopf–Cole y esquemas numéricos}} \\[0.4em]
	\large Informe de anteproyecto y estado de avance}
\author{%
	{Anthonny Flores Rojas}\thanks{{Escuela de Matemáticas, Universidad de Costa Rica}. \texttt{{anthonny.flores@ucr.ac.cr}}} \and
	{Leonardo Vega Aragón}\thanks{{Escuela de Matemáticas, Universidad de Costa Rica}. \texttt{\ph{@ucr.ac.cr}}} \and
	{Melody Campos Pérez}\thanks{{Escuela de Matemáticas, Universidad de Costa Rica}. \texttt{\ph{@ucr.ac.cr}}} \and
	{Sebastián Calderón Segura}\thanks{{Escuela de Matemáticas, Universidad de Costa Rica}. \texttt{\ph{@ucr.ac.cr}}}}
\date{Curso: {Análisis Numérico (MA0501)} \\
	Profesor: {Esteban Segura} \\[0.3em] \today}

\begin{document}
	\maketitle
	
	% ---------------------------------------------------------------------
	\begin{abstract}
		\noindent
		\ph{Resumen del proyecto en 5 a 8 líneas: qué problema se estudia, qué se ha
			demostrado o logrado hasta ahora, qué método se usa y qué falta por hacer.}
	\end{abstract}
	
	\medskip
	\noindent\textbf{Clasificación MSC 2020:} \ph{65Mxx}
	(primaria); \ph{35Bxx} (secundaria).\\
	\noindent\textbf{Palabras clave:} \ph{palabra clave 1, palabra clave 2, palabra clave 3}.
	
	% ---------------------------------------------------------------------
	\section{Introducción}\label{sec:intro}
	
	\ph{Presentar el problema: qué se estudia, por qué es interesante y qué se
		sabe hasta ahora. Incluir uno o dos ejemplos sencillos que motiven la
		pregunta.}
	
	\paragraph{Resultados principales.}
	\ph{Resumir qué se demuestra o se obtiene en este proyecto, citando la
		bibliografía de base} \cite{Ref1,Ref2}\ph{.}
	
	\begin{theorem}[\ph{Nombre del teorema}]\label{thm:principal1}
		\ph{Enunciado preciso del primer resultado principal.}
	\end{theorem}
	
	\begin{theorem}[\ph{Nombre del teorema}]\label{thm:principal2}
		\ph{Enunciado preciso del segundo resultado principal.}
	\end{theorem}
	
	\paragraph{Organización.}
	La \cref{sec:prelim} reúne los preliminares. En la \cref{sec:principal} se
	demuestran los resultados principales. La \cref{sec:experimentos} presenta
	la verificación computacional. La \cref{sec:estado} resume el estado del
	proyecto y el trabajo restante. \ph{Ajustar esta descripción a las secciones
		que realmente se usen.}
	
	% ---------------------------------------------------------------------
	\section{Preliminares}\label{sec:prelim}
	
	\subsection{Definiciones y notación}
	
	\ph{Introducir los objetos básicos y la notación que se usará en todo el
		documento.}
	
	\begin{definition}\label{def:principal}
		\ph{Definición del concepto principal.}
	\end{definition}
	
	\begin{example}
		\ph{Ejemplo sencillo que ilustre la definición anterior.}
	\end{example}
	
	\subsection{Resultados previos}
	
	\begin{lemma}\label{lem:previo}
		\ph{Enunciado de un resultado previo que se necesitará más adelante.}
	\end{lemma}
	
	\begin{proof}
		\ph{Demostración del lema.}
	\end{proof}
	
	\begin{proposition}\label{prop:previa}
		\ph{Enunciado de otra proposición auxiliar.}
	\end{proposition}
	
	\begin{proof}
		\ph{Demostración de la proposición.}
	\end{proof}
	
	% ---------------------------------------------------------------------
	\section{Resultados principales}\label{sec:principal}
	
	\begin{proof}[Demostración del \cref{thm:principal1}]
		\ph{Desarrollo completo de la demostración, indicando en cada paso qué
			resultado previo se usa, por ejemplo el \cref{lem:previo}.}
	\end{proof}
	
	\begin{lemma}\label{lem:auxiliar}
		\ph{Enunciado del lema auxiliar necesario para el segundo teorema.}
	\end{lemma}
	
	\begin{proof}
		\ph{Demostración del lema auxiliar.}
	\end{proof}
	
	\begin{proof}[Demostración del \cref{thm:principal2}]
		\ph{Desarrollo completo de la demostración, usando el \cref{lem:auxiliar}.}
	\end{proof}
	
	\begin{remark}
		\ph{Comentario sobre el resultado: limitaciones, casos particulares,
			posibles generalizaciones.}
	\end{remark}
	
	% ---------------------------------------------------------------------
	\section{Experimentos computacionales}\label{sec:experimentos}
	
	\ph{Explicar qué se calculó o verificó con computadora, con qué herramienta
		y para qué rango de valores. Si el proyecto no incluye cálculos, eliminar
		esta sección.}
	
	\begin{lstlisting}
		# Escribir aqui el codigo usado (Python, SageMath, etc.)
	\end{lstlisting}
	
	\ph{Explicar los resultados obtenidos y qué conclusión se saca de ellos, en
		relación con la \cref{tab:resultados}.}
	
	\begin{table}[ht]
		\centering
		\renewcommand{\arraystretch}{1.2}
		\begin{tabular}{@{}ccc@{}}
			\toprule
			\ph{Columna 1} & \ph{Columna 2} & \ph{Columna 3} \\
			\midrule
			\ph{dato} & \ph{dato} & \ph{dato} \\
			\ph{dato} & \ph{dato} & \ph{dato} \\
			\ph{dato} & \ph{dato} & \ph{dato} \\
			\bottomrule
		\end{tabular}
		\caption{\ph{Descripción de los datos mostrados.}}\label{tab:resultados}
	\end{table}
	
	% ---------------------------------------------------------------------
	\section{Estado del proyecto y trabajo restante}\label{sec:estado}
	
	\subsection{Resumen del avance}
	
	\ph{Introducir brevemente la tabla: a qué fecha corresponde y qué se ha
		logrado.} Los estados posibles son \done, \wip\ y \todo. La
	\cref{tab:estado} indica el responsable de cada objetivo.
	
	\begin{table}[ht]
		\centering
		\renewcommand{\arraystretch}{1.25}
		\begin{tabular}{@{}l p{5.6cm} l c@{}}
			\toprule
			\textbf{Objetivo} & \textbf{Descripción} & \textbf{Responsable} & \textbf{Estado} \\
			\midrule
			O1 & \ph{Descripción del objetivo 1} & \ph{Nombre} & \done \\
			O2 & \ph{Descripción del objetivo 2} & \ph{Nombre} & \wip \\
			O3 & \ph{Descripción del objetivo 3} & \ph{Nombre} & \wip \\
			O4 & \ph{Descripción del objetivo 4} & \ph{Nombre} & \todo \\
			O5 & \ph{Redacción del informe final} & \ph{Todos} & \todo \\
			\bottomrule
		\end{tabular}
		\caption{Estado de los objetivos a la fecha del informe.}\label{tab:estado}
	\end{table}
	
	\subsection{Trabajo en curso}
	
	\ph{Describir en qué se está trabajando ahora: qué se quiere demostrar o
		construir y hasta dónde se ha llegado.}
	\begin{equation}\label{eq:pendiente}
		\ph{fórmula o resultado que se quiere probar}
	\end{equation}
	
	\begin{conjecture}\label{conj:pendiente}
		\ph{Conjetura o resultado pendiente de demostrar (opcional).}
	\end{conjecture}
	
	\subsection{Dificultades y plan}
	
	\begin{itemize}[leftmargin=2em]
		\item \textbf{Dificultad:} \ph{descripción del problema encontrado.}
		\emph{Plan:} \ph{cómo se piensa resolver.}
		\item \textbf{Dificultad:} \ph{descripción del problema encontrado.}
		\emph{Plan:} \ph{cómo se piensa resolver.}
	\end{itemize}
	
	\subsection{Cronograma}
	
	\begin{enumerate}[leftmargin=2em]
		\item \ph{Actividad 1} \hfill \emph{(\ph{semanas 1--2})}
		\item \ph{Actividad 2} \hfill \emph{(\ph{semanas 2--4})}
		\item \ph{Redacción del informe final y presentación} \hfill \emph{(\ph{semanas 4--6})}
	\end{enumerate}
	
	% ---------------------------------------------------------------------
	\begin{thebibliography}{9}
		
		\bibitem{Lindgren2020} J. Lindgren, ``Efficient Markets and Contingent Claims Valuation: An Information Theoretic Approach,'' \emph{Entropy}, vol. 22, no. 11, art. 1283, Nov. 2020, doi: 10.3390/e22111283.
		
		\bibitem{Ref2}
		\ph{Autor(es), \emph{Título del artículo}, revista \textbf{volumen} (año), páginas.}
		
	\end{thebibliography}
	
\end{document}